\documentclass[%
  aps, amsmath, amssymb, onecolumn, 11pt
]{revtex4-2}

\usepackage{graphicx}
\graphicspath{{figures/}}
\usepackage{orcidlink}
\usepackage{mathtools}
\usepackage{booktabs}

% -------------- arXiv-Style Margin Stamp ----------------
\usepackage[style=iso]{datetime2}
\usepackage{FiraMono}
\usepackage{tikz}

\AddToHook{shipout/background}{
  \begin{tikzpicture}[remember picture, overlay]
    \node [rotate=-90, color=gray!80, font=\fontfamily{FiraMono-TLF}\selectfont\footnotesize] at ([xshift=-1.2cm]current page.east) {WORKING DRAFT \quad $\bullet$ \quad SUPPLEMENTAL MATERIAL \quad $\bullet$ \quad \DTMnow \quad $\bullet$ \quad doi.org/10.17605/OSF.IO/EV8H6 \quad $\bullet$ \quad Brian S. Mulloy \quad $\bullet$ \quad bmulloy@umich.edu};
  \end{tikzpicture}
}
% --------------------------------------------------------

% --- SM-specific counter renumbering ---
\renewcommand{\thesection}{S\arabic{section}}
\renewcommand{\thesubsection}{S\arabic{section}.\arabic{subsection}}
\renewcommand{\thefigure}{S\arabic{figure}}
\renewcommand{\thetable}{S\arabic{table}}
\renewcommand{\theequation}{S\arabic{equation}}

\begin{document}

\title{Supplemental Material for\\``Wigner Resolution as the Reciprocal of Phase-Space Extent''}

\author{Brian S. Mulloy\orcidlink{0000-0002-1803-3172}}
\email{bmulloy@umich.edu}

\date{\today}

\maketitle

% ======================================================================
% S1 --- Notation and conventions
% ======================================================================

\section{Notation and conventions}
\label{sec:notation}

% [To be drafted.]
% Topics:
% - Fermi blob: $A = \pi \Delta q \Delta p \geq h/2$
% - Quantum blob at angle $\theta$: $a_\theta = \pi \delta q_\theta \delta p_\theta = h/2$
% - Principal-axis cases: $a_q \equiv a_{\theta=0}$, $a_p \equiv a_{\theta=\pi/2}$
% - Polar-dual identity: $\delta q_\theta \, \delta p_\theta = \hbar$
% - Inscription bounds: $\delta q \leq \Delta q$, $\delta p \leq \Delta p$
% - Convention: $\hbar$ for dispersions and Wigner kernel; $h$ for phase-space areas.
%   Mention only if needed; otherwise allow the reader to absorb from context.

% ======================================================================
% S2 --- Full action-units table
% ======================================================================

\section{Action-units table for phase-space methods}
\label{sec:fulltable}

% [To be drafted.]
% The full 16-row table from the body's truncated version.
% Columns: Model | Native subfield | Kernel action structure | Map to $\delta q = \hbar/\Delta p$ | Star ranking
% Brief paragraph framing before the table.

% ======================================================================
% S3 --- Effective resolution derivations
% ======================================================================

\section{Effective resolution: derivations}
\label{sec:derivations}

\subsection{Generalized Husimi $Q_{\lambda\theta}$}
\label{sec:genhusimi}

% [To be drafted.]
% Show: kernel $\lambda \cdot 1/\lambda = \hbar$; setting $\lambda = \hbar/\Delta p$ recovers $P_{\delta q}$.
% Cite Wódkiewicz, Appleby, Wünsche-Bužek.

\subsection{Truncated Wigner Approximation}
\label{sec:twa}

% [To be drafted.]
% Show: TWA samples at $\Delta q$. At saturation $\Delta q \cdot \Delta p \approx \hbar$,
% the sampling width equals the polar-dual scale $\delta q = \hbar/\Delta p$.
% The silhouette argument.

\subsection{Cahill--Glauber $W^{(s)}$}
\label{sec:cahillglauber}

% [To be drafted.]
% Show: $a_s = \pi(1-s)\hbar/2$. State-matched $s$ value.

\subsection{Husimi $Q$}
\label{sec:husimi}

% [To be drafted.]
% Show: fixed kernel area $h/2$; matches $\delta q$ only at vacuum.

\subsection{Maximum-likelihood reconstruction with Fock-state truncation}
\label{sec:maxlik}

% [To be drafted.]
% Show: Fock truncation $N_{\max}$; resolution per cell $\sim \hbar/\Delta p$ for state with
% covariance $\Delta p^2 \sim N_{\max}\hbar$. Polar-dual scale emerges from action-counting.

% ======================================================================
% S4 --- Notes on remaining methods
% ======================================================================

\section{Notes on remaining methods}
\label{sec:remaining}

% [To be drafted.]
% One paragraph each:
% - Herman--Kluk frozen Gaussian
% - Heller frozen/thawed Gaussian
% - Pattern-function tomography
% - Positive-P representation
% - Cohen-class Born--Jordan
% - Semidefinite programming tomography
% - Bayesian / MCMC tomography
% - Compressed sensing tomography
% - Husimi-based open quantum systems

% ======================================================================
% S5 --- Computational details
% ======================================================================

\section{Computational details}
\label{sec:computational}

% [To be drafted.]
% - R pipeline overview
% - Eight-state list with covariance parameters
% - Numerical grid choices for Wigner-function evaluation
% - OSF deposit pointer

\bibliography{main}

\end{document}
